\chapter{Conclusion and Future Work}
\label{ch:conclusion}

\section{Conclusion}

This project set out to build a complete and honest fraud detection system on a severely
imbalanced dataset, and to work out how much confidence its results actually deserve.

On the engineering objectives it succeeded. The pipeline de-duplicates before splitting and
keeps all resampling inside the training fold. Eight model pipelines are trained and
compared under a protocol in which the test set influences nothing but the final report. The
best model reaches a PR-AUC of 0.828, catching 70 of 95 frauds at two false alarms. The
whole system is served through an API and an interactive web application in which the
decision threshold is a live control rather than a hidden constant.

On the statistical objectives it produced a more interesting result than we expected. Once
we applied confidence intervals, paired significance tests, and cross-validation, the
apparent ordering of the leading models did not hold up. The champion's bootstrap interval
spans $[0.751, 0.899]$. None of the 28 pairwise McNemar tests is significant after Holm
correction. Cross-validation shows fold-to-fold variation several times larger than the
differences between models, and it reverses the ranking. The top four models cannot be told
apart on this dataset.

The differences that are real lie elsewhere. Resampling strategy produces a spread thirteen
times wider than algorithm choice, with random undersampling collapsing to 0.548 because
balancing the classes destroys 99.7 percent of the training data. Calibration separates the
model families by two orders of magnitude. And the learning curve, still rising at the full
training set, points to the real bottleneck, which is not the algorithm but the number of
fraud examples available to learn from.

The main lesson of the project is therefore about method. On a dataset with 473 positive
cases, a single train/test split cannot support the model-ranking claims that are routinely
made from it, and the extra work needed to discover that, being intervals, paired tests and
resampled folds, is small compared to the false confidence it prevents.

\section{Summary of Contributions}

\begin{itemize}
  \item A reproducible, leakage-controlled pipeline with de-duplication before splitting and
        resampling confined to the training fold.
  \item A comparison of eight pipelines in which every headline figure carries a confidence
        interval and every pairwise difference is tested for significance, with the
        conclusion drawn being the one the evidence supports rather than the one the point
        estimates suggest.
  \item A controlled resampling study that separates the sampling step from the algorithm,
        with a concrete explanation for why undersampling fails.
  \item A calibration, cost, lift, feature-importance, and learning-curve analysis of the
        deployed model, including the finding that the simple error-count threshold is
        within 0.3 percent of cost-optimal under an amount-weighted cost model.
  \item A served system with a model registry, exposing model choice and decision threshold
        as live controls, together with a three-dimensional rendering of the full decision
        surface.
  \item A reporting chain in which every table is generated from the analysis outputs, so
        the document cannot drift away from the code.
\end{itemize}

\section{Future Work}

Nested cross-validation with tuning inside the loop is the clearest methodological
improvement available. Tuning every algorithm equally, inside each fold,
would remove the asymmetry identified in Section~\ref{sec:notestablish} and give a fair
estimate of what each algorithm can reach, rather than comparing one tuned model against
three untuned ones.

Collecting more labelled fraud is the next priority. The learning curve has not levelled
off, so this is worth more than further algorithmic work.

Local explanations would also be valuable. \gls{shap} values \citep{lundberg2017unified}
would show which
features pushed an individual transaction toward its score. With anonymised components this
cannot produce a human-readable reason, but it would support consistency auditing and help
meet the explainability requirements that apply to automated decisions in financial
services.

Sequence models are a natural extension. Given per-card histories with genuine timestamps, recurrent
architectures such as \gls{lstm} \citep{hochreiter1997long} could capture patterns that a
single-transaction classifier cannot see. A burst of small test purchases is the standard
example, and the small median fraud amount in this dataset hints at exactly that behaviour.

A calibrated cost-sensitive deployment is possible too. Combining the well-calibrated
probabilities of Section~\ref{sec:calibration} with a measured cost model would let the
threshold be set by expected loss directly, and adjusted automatically as review capacity
changes.

Finally, a production deployment would need drift monitoring and scheduled retraining, so
that distribution shift is detected and the model is refitted on a schedule. A two-day
training window says nothing about stability over time.
